
\section{La competencia entre crecimiento de grieta y plasticidad}
La intensidad elástica singular describe un campo intermedio: suficientemente lejos del proceso microscópico de la punta y suficientemente cerca para que domine el término de intensidad. Una estimación de escala de zona plástica en tensión plana es
\begin{equation}r_p\approx\frac{1}{2\pi}\left(\frac{K_I}{\sigma_y}\right)^2.\end{equation}
En una idealización de deformación plana se utiliza un factor diferente, frecuentemente $1/(6\pi)$. Son estimaciones de orden de magnitud. No sustituyen las condiciones completas de un ensayo ni permiten decidir por sí solas que la mecánica lineal es válida. Deben compararse con grieta, ligamento, espesor y gradientes.

Si la plasticidad deja de ser pequeña, energía liberada y deformación plástica se acoplan. Parámetros como la integral $J$ y la apertura de punta pueden ser más adecuados, bajo sus propias condiciones. Una curva de resistencia al crecimiento puede mostrar que la resistencia no permanece constante durante extensión estable. Un valor de $K_{IC}$ no debe usarse como una propiedad universal que cubra cualquier geometría, régimen y mecanismo.
\section{Qué información lleva realmente un factor geométrico}
La expresión $Y\sigma\sqrt{\pi a}$ separa una escala de tensión, una escala de tamaño y una corrección de geometría y carga. $Y$ puede depender de $a/W$, profundidad, relación de aspecto, forma de la superficie y distribución de tensión. Una grieta de borde y una grieta central no comparten automáticamente definición de $a$ ni factor.

Si $Y=Y(a)$, despejar $a_c$ con un valor fijo puede ser inconsistente. Se formula $f(a)=Y(a)\sigma\sqrt{\pi a}-K_c$ y se busca una raíz dentro del dominio físico. La comprobación posterior sustituye la raíz en la ecuación original y verifica ligamento y otras hipótesis. El problema P10 del cuaderno ilustra esta diferencia sin presentar su solución geométrica como válida para toda pieza.
\section{Del defecto observado a una decisión documentada}
La medida de un defecto incluye incertidumbre, orientación, acceso y método. Una indicación superficial puede requerir caracterización adicional; una longitud proyectada no determina siempre profundidad. El análisis usa una idealización explícita de defecto y justifica cómo se relaciona con la evidencia. Después considera las cargas máximas, el crecimiento entre inspecciones y la probabilidad de detección.

La tolerancia al daño integra resistencia residual, evolución del defecto y capacidad de inspección. Si una inspección no cubre una zona crítica, no se compensa solo reduciendo un tamaño inicial en la fórmula. La documentación aplicable determina la actuación real; el cálculo académico enseña a detectar qué hipótesis faltan antes de aceptar una conclusión.
\repaso{Si dos piezas tienen el mismo $K_I$, ¿tienen necesariamente igual seguridad?}{No. Pueden diferir tenacidad aplicable, espesor, ambiente, plasticidad, historia de carga, tensiones residuales y capacidad de detección. La igualdad del parámetro aplicado no iguala todos los mecanismos ni la evidencia.}
